\scanpage{50}
% archive foliation: 49
\hfill - 36 -\hfill n° 395

\deleted{
que $\|u_n\|<C$, $x_n=\sum_{k=1}^{n}\gamma(1-\gamma)^k u_k+(1-\gamma)^n e$ admette un inverse $t_n$, et $\|t_nz-t_{n-1}z\|<\dfrac{\delta}{2^n}$, $\|t_1z-z\|<\dfrac{\delta}{2}$. Posant $y_n=t_nz$, les $y_n$ tendent vers une limite $y\in A$, $x_n$ tend vers $x=\sum_{k=1}^{\infty}\gamma(1-\gamma)^{k-1}u_k$; x et y ont les propriétés requises).

23) a) Soient E un espace de Banach, $f:\mathbb{C}\longrightarrow E$ une fonction entière, et $\vartheta\longrightarrow p(\vartheta)=\sum_{\mu=k}^{m}c_\mu e^{\mu i\vartheta}$ un polynôme trigonométrique, $(c_k\ne0)$. On suppose que $|p(\vartheta)|\,\|f(re^{i\vartheta})\|\leq Mr^\alpha$ pour $r\geq r_0$, (M, $r_0$, $\alpha$ étant des constantes $\geq0$). Alors f est un polynôme de degré $\leq\alpha$. (Soient $f(\zeta)=\sum_{\nu=0}^{\infty}a_\nu\zeta^\nu$, et $J(r,p)=r^{-p}\int_0^{2\pi}p(\vartheta)f(re^{i\vartheta})e^{-(k+p)i\vartheta}\,d\vartheta$. On a $J(r,p)\longrightarrow0$ quand $r\longrightarrow+\infty$ si $p>\alpha$, et d'autre part $\lim J(r,p)=c_ka_p$).

b) Prouver de même que, si $\alpha$ est entier et si $r^{-\alpha}|p(\vartheta)|\,\|f(re^{i\vartheta})\|\longrightarrow0$ quand $r\longrightarrow+\infty$, f est un polynôme de degré $<\alpha$.

c) Soient A une algèbre de Banach \unsure{unipollente|unipotente}, x un élément quasi-nilpotent de A. On suppose que, pour un certain entier $n>0$, on a $r^{-n}|\cos^n\vartheta|\,\|(e-\lambda\unsure{x|z})^{-1}\|\longrightarrow0$ quand $r\longrightarrow+\infty$, $(\lambda=re^{i\vartheta})$. Alors $x^n=0$. (Utiliser b)).

\P\quad 24) a) Soient E un espace de Banach, $x:\mathbb{C}-\{1\}\longrightarrow E$ une fonction entière de $1/(\zeta-1)$, $x(\zeta)=\sum_{n=0}^{\infty}a_n\zeta^n$, $x(\zeta)=\sum_{n=0}^{\infty}b_n\zeta^{-n}$ les développements de x pour $|\zeta|<1$, $|\zeta|>1$ respectivement. On suppose $\|a_n\|=\unsure{o|O}(n^\alpha)$, $\|b_n\|=\unsure{o|O}(n^\alpha)$, $(\alpha>0)$. Alors $x(\zeta)$ est un polynôme en $1/(\zeta-1)$ de degré $<\alpha+1$. (Soit $\varepsilon>0$. On a $\|x(\zeta)\|\leq\varepsilon(1-r)^{-1-\alpha}$ pour $r_\varepsilon\leq r=|\zeta|<1$, $\|x(\zeta)\|\leq\varepsilon(1-r^{-1})^{-1-\alpha}$ pour $1<r<r'_\varepsilon$. Posant $1-\zeta=(\omega+\frac12)^{-1}$, $\omega=Re^{i\vartheta}$ montrer que, pour $0\leq r<1$ et $R\geq1$, on a $(1-r)^{-1}\leq\dfrac94\dfrac{R}{\cos\vartheta}$ et que, pour $1<r$ et $R\geq1$, $(1-r^{-1})^{-1}\leq\dfrac94\dfrac{R}{\cos\vartheta}$. Poser $y(\omega)=x(\zeta)$ et montrer que
}

